\question[5]
In einem Programm sollen folgende Daten gespeichert werden. Gib jeweils einen geeigneten Python-Datentyp
(\texttt{int}, \texttt{float}, \texttt{bool} oder \texttt{str}) an und begründe kurz.
\begin{parts}
\part die Anzahl der Schüler eines Kurses
\part die Körpergröße in Metern
\part ob ein Schüler volljährig ist
\part ein Kontostand, mit dem exakt gerechnet werden soll
\part die Postleitzahl, z.B. 01067
\end{parts}
In Java soll die Einwohnerzahl der Erde (etwa 8 Milliarden) gespeichert werden.
Zur Erinnerung: In Java speichert \texttt{int} eine ganze Zahl mit 32 Bit und \texttt{long} mit 64 Bit,
beide im Zweierkomplement. Welcher der beiden Typen ist nötig? Begründe mit einer Abschätzung des Wertebereichs.

\begin{solutionbox}{7cm}
\begin{parts}
\part \texttt{int}: eine Anzahl ist immer ganzzahlig.
\part \texttt{float}: eine Größe hat Nachkommastellen.
\part \texttt{bool}: es gibt nur die zwei Möglichkeiten ja und nein.
\part \texttt{int}, und zwar in Cent: Mit \texttt{float} entstehen Rundungsfehler, mit ganzen Zahlen nicht.
\part \texttt{str}: Mit einer Postleitzahl wird nicht gerechnet, und als \texttt{int} ginge die führende 0 verloren.
\end{parts}
\texttt{long}: Der größte \texttt{int}-Wert ist $2^{31} - 1$. Mit $2^{10} \approx 10^3$ ist
$2^{31} = 2 \cdot (2^{10})^3 \approx 2 \cdot 10^9$, also nur etwa 2 Milliarden. Das reicht nicht.
Bei \texttt{long} ist $2^{63} = 8 \cdot (2^{10})^6 \approx 8 \cdot 10^{18}$, das reicht.
\end{solutionbox}
